\subsection{Klíčové funkce systému}
\label{subsec:kelvin-features}

Systém Kelvin nabízí několik funkcí, které jsou podstatné jak pro běžný provoz, tak pro plánovanou integraci s LLM\@.

\begin{itemize}
    \item \textbf{Automatické evaluace} -- po každém odevzdání systém automaticky zkompiluje studentův kód, spustí jej proti sadě předpřipravených testovacích vstupů a porovná výstupy s očekávanými hodnotami.
    Zdrojový kód je tak testován na jednotném prostředí, což zajišťuje spravedlivé hodnocení funkční správnosti.
    Výsledky jednotlivých testů jsou studentovi zobrazeny okamžitě, což umožňuje rychlou iteraci nad řešením bez nutnosti čekat na zpětnou vazbu od vyučujícího.

    \item \textbf{Automatické udělování bodů} -- na základě výsledků automatických testů může systém přidělit body automaticky podle poměru úspěšných testů.
    Vyučující tak nemusí manuálně hodnotit funkční správnost a může se soustředit na kvalitativní aspekty řešení, jako je čitelnost kódu, vhodnost zvolené datové struktury nebo efektivita algoritmu.

    \item \textbf{Řádkové a celkové komentáře} -- vyučující může označit libovolný řádek zdrojového kódu a připsat k němu připomínku.
    Kromě řádkových komentářů systém umožňuje i vložení celkového komentáře k odevzdání, který není vázán na konkrétní řádek a slouží jako celkové shrnutí zpětné vazby.
    Komentáře jsou studentovi zobrazeny přímo v kontextu kódu, takže okamžitě vidí, ke které konkrétní části řešení se zpětná vazba vztahuje.
    Tato vlastnost je z pohledu integrace LLM zcela zásadní, protože LLM generuje komentáře přirozeně ve formátu \uv{k řádku X patří připomínka Y}.
    Výstup modelu tak lze přímo mapovat na existující datový model systému bez nutnosti zavádět nové typy do databáze.
\end{itemize}

\endinput
